\appendixexercisegroup

\begin{enumerate}
\def\labelenumi{\arabic{enumi}.}
\item
  设S是一组符号，设A是 \(\mathbf{L}_{0}(\mathbf{S})\) 中的一个词，且 B、C 为 A 的两个显著段。则要么 B 是 C 的段，要么 C 是 B 的段，要么 B 与 C 不相交。
\item
  设 S 为一个符号集，设 A 为 \(\mathrm{L}_{0}(\mathrm{S})\) 中的一个显著词，并假设 \(A = A'BA''\) ，其中 B 是显著的。证明：若 C 是一个显著词，则词 \(A'CA''\) 是显著的。（使用练习1。）
\item
  令 E 为一个集合，并令 f 为 E 的一个映射。 \(E \times E\) 进入E（“内部合成律”）。设S为一组符号，它是E与仅含单个元素f的集合的不交并。～令 \(n(f) = 1\) ，以及 \(n(x) = 0\) 对于所有 \(x \in E\) .
\end{enumerate}

\begin{enumerate}
\def\labelenumi{(\alph{enumi})}
\item
  设 M 为 \(L_{0}(S)\) 的显著词组成的集合。证明存在唯一的映射 v 将 M 映入 E，且满足以下条件：(1) \(v(x) = x\) 对于所有 \(x \in E\) 成立；(2) 若 A 和 B 是两个显著词，则 \(v(fAB) = f(v(A), v(B))\) .
\item
  对于每个词 \(A = (s_i)_{0 \leqslant i \leqslant n}\) （ \(\mathbf{L}_0(\mathbf{S})\) 中的），令 \(A^*\) 为词 \((s_{i_k})\) ，其中 \(i_k\) 是指标 \(i\) 使得 \(s_i \neq f\) ，按升序排列。两个词 \(A, B\) （ \(\mathbf{L}_0(\mathbf{S})\) 中的）被称为相似的，如果 \(A^* = B^*\) 。证明如果组合法则 \(f\) 是结合的，即如果
\end{enumerate}

\[
f (f (x, y), z) = f (x, f (y, z)),
\]

那么， \(v(A) = v(B)\) ，只要 \(A\) 和 \(B\) 是相似的显著词（“结合性的一般定理”）。（一个显著词

\[
A = (s _ {i}) _ {0 \leqslant i \leqslant 2 n}
\]

被称为正规的，如果 \(s_{i}=f\) 为了 \(i=0,2,4,\ldots,2n-2\) ，以及 \(s_{i}\neq f\) 对于其他指标。证明每个显著词都相似于唯一的正规词 \(A'\) ，并证明 \(v(A)=v(A')\) 通过对A的长度进行归纳。）

\begin{enumerate}
\def\labelenumi{\arabic{enumi}.}
\setcounter{enumi}{3}
\tightlist
\item
  设 \(A\) 是理论 \(\mathfrak{G}\) 中的一个项或一个关系。考虑如下定义的组装序列。首先，写出 \(A\) 。如果 \(A\) 由单个字母组成，则在此停止。否则，接下来写出A的前驱组装或前驱组装们（如果A以 \(\tau\) 开头，则任意选择其中一个前驱组装）。然后写出前述那些不是由单个字母组成的组装的前驱组装或前驱组装们；依此类推。
\end{enumerate}

\begin{enumerate}
\def\labelenumi{(\alph{enumi})}
\item
  如果我们颠倒这个组装序列的顺序，就得到一个形成性构造。
\item
  设B是A的一个平衡段，且B中没有符号在A中与B外部的符号相连接。证明B要么是一个项，要么是一个关系（使用（a）和练习1）。
\item
  若B是项（resp. 关系），则将A中的B替换为项（resp. 关系）。证明：若A是项，则如此得到的组装是项；若A是关系，则如此得到的组装是关系。
\end{enumerate}

\begin{enumerate}
\def\labelenumi{\arabic{enumi}.}
\setcounter{enumi}{4}
\item
  设A是理论C中的一个组装，设T是G中的一个项，设x是一个字母。如果 \((T|x)A\) 是一个项（相应地，一个关系），则 A 是一个项（相应地，一个关系）。（使用练习 4。）
\item
  一个理论中的关系 \(\mathcal{G}\) 被称为逻辑上不可约的，如果它以关系符号开头。设 \(R_{1}, R_{2}, \ldots, R_{n}\) 是逻辑上不可约的关系，在 \(\mathcal{G}\) . 一种基于逻辑的构造 \(R_{1}, R_{2}, \ldots, R_{n}\) 任何序列 \(A_{1}, A_{2}, \ldots, A_{p}\) 组件组装中的 \(\mathcal{G}\) 使得每个 \(A_{i}\) 满足以下条件之一：(1) \(A_{i}\) 是以下关系之一 \(R_{1}, R_{2}, \ldots, R_{n}\) ； (2) 存在一个组装 \(A_{j}\) 在前 \(A_{i}\) 使得 \(A_{i}\) 是 \(\neg A_{j}\) ； (3) 存在两个组装 \(A_{j}\) 并且 \(A_{k}\) 在前 \(A_{i}\) 使得 \(A_{i}\) 是 \(\vee A_{j}A_{k}\) .
\end{enumerate}

\begin{enumerate}
\def\labelenumi{(\alph{enumi})}
\tightlist
\item
  证明以
\end{enumerate}

\[
R _ {1}, R _ {2}, \dots , R _ {n}
\]

这些规则是准则 CS8、CS9（第 1 节）、CS5（§1，第 2 节）和 CS6（§3，第 4 节）的直接推论。 \(\mathfrak{C}\) 。出现在这种逻辑构造中的关系被称为由 \(R_{1}, R_{2}, \ldots, R_{n}\) .

\begin{enumerate}
\def\labelenumi{(\alph{enumi})}
\setcounter{enumi}{1}
\item
  如果R和S是由 \(R_{1}, R_{2}, \ldots, R_{n}\) 逻辑构造的，那么 \(\neg R, \vee RS, \Longrightarrow RS\) ，“R且 \(S\)'' , \(R \Longleftrightarrow S\) .
\item
  设R是C中的一个关系。考虑如下定义的关系序列。首先，写出R。如果R是逻辑不可约的，则在此停止。如果不是，则接着写出R的前驱组装（这些是确定的关系）。然后写出那些不是逻辑不可约的前驱组装的前驱组装；依此类推。设 \(R_{1}, R_{2}, \ldots, R_{n}\) 是这一构造中产生的不同的逻辑不可约关系；它们被称为 R 的逻辑分量。证明 R 是由其逻辑分量逻辑构造的，但如果序列中的某个关系
\end{enumerate}

\(R_{1}, R_{2}, \ldots, R_{n}\) 被省略，则 R 不能由其余关系逻辑构造 \((n - 1)\) 关系。

\begin{enumerate}
\def\labelenumi{(\alph{enumi})}
\setcounter{enumi}{3}
\tightlist
\item
  设 R 是一个关系，且设 \(R_{1}, R_{2}, \ldots, R_{n}\) 是互不相同的逻辑不可约关系，使得（1）R 由 \(R_{1}, R_{2}, \ldots, R_{n}\) 逻辑构造而成，（2）如果我们从序列 \(R_{1}, R_{2}, \ldots, R_{n}\) 中省略一个关系，则 R 不能由剩余关系逻辑构造而成。证明 \(R_{1}, R_{2}, \ldots, R_{n}\) 是 R 的逻辑分量。
\end{enumerate}

\begin{enumerate}
\def\labelenumi{\arabic{enumi}.}
\setcounter{enumi}{6}
\tightlist
\item
  设 \(R_{1}, R_{2}, \ldots, R_{n}\) 是理论 C 中互不相同的逻辑不可约关系（习题 6）。设 \(A_{1}, A_{2}, \ldots, A_{m}\) 是一个以 \(R_{1}, R_{2}, \ldots, R_{n}\) 为基底的逻辑构造。假设每个 \(R_{j}\) 被赋予符号0、1中的一个。然后我们按以下规则给每个 \(A_{i}\) 赋予符号0、1中的一个： (1) 若 \(A_{i}\) 与 \(R_{j}\) 相同，给 \(A_{i}\) 与 \(R_{j}\) 相同的符号； (2) 若 \(A_{i}\) 与 \(\neg A_{j}\) 相同，其中 \(A_{j}\) 先于 \(A_{i}\) ，则赋予 \(A_{i}\) 符号1（相应地，0），如果 \(A_{j}\) 具有符号0（相应地，1）；（3）如果 \(A_{i}\) 与 \(\vee A_{j}A_{k}\) 相同，给出 \(A_{i}\) 符号1，如果 \(A_{j}\) 和 \(A_{k}\) 两者符号均为1；否则给出 \(A_{i}\) 符号0。（我们正在应用以下“符号规则”：
\end{enumerate}

\[
\neg 0 = 1, \quad \neg 1 = 0, \quad \vee 1 1 = 1, \quad \vee 1 0 = \vee 0 1 = \vee 0 0 = 0.)
\]

\begin{enumerate}
\def\labelenumi{(\alph{enumi})}
\item
  证明只有一种方式为每个 \(A_{i}\) 根据这些规则赋予符号。
\item
  如果 R 是由 \(R_{1}\) , \(R_{2}\) , \(\ldots\) , \(R_{n}\) ，那么R所携带的符号与逻辑构造无关，以
\end{enumerate}

\[
R _ {1}, R _ {2}, \dots , R _ {n},
\]

中的一个证明，在其中 \(\pmb{R}\) 出现。

\begin{enumerate}
\def\labelenumi{(\alph{enumi})}
\setcounter{enumi}{2}
\item
  如果R和S是由 \(R_{1}, R_{2}, \ldots, R_{n}\) ，且若R所携带的符号 \(\Longrightarrow RS\) 为0，则S所携带的符号为0。
\item
  从现在起，假设 \(\mathfrak{C}\) 的公理仅由方案S1至S4提供。设 \(\pmb{R}\) 是 \(\mathfrak{C}\) 中的一个定理，并设 \(\pmb{R}_1, \pmb{R}_2, \ldots, \pmb{R}_n\) 是其逻辑组成部分。然而，试表明符号0和1被赋予 \(\pmb{R}_1, \pmb{R}_2, \ldots, \pmb{R}_n\) ，且 \(\pmb{R}\) 所携带的符号为 0。（首先对 \(\pmb{R}\) 是 \(\mathfrak{C}\) 的公理的情形加以证明；在一般情形，考虑 \(\pmb{R}\) 的一个证明并应用 (c)) \(\mathfrak{C}\) .
\item
  设R为C中一个逻辑不可约关系。证明R和“非R”都不是C中的定理。特别地，C是无矛盾的（使用(d)）。
\item
  设 \(\mathbf{R}_1, \mathbf{R}_2, \ldots, \mathbf{R}_n\) 是 \(\mathfrak{C}\) 中互不相同的逻辑不可约关系。考虑所有形如“ \(\mathbf{R}_1'\) 要么 \(\mathbf{R}_2'\) 要么 \ldots{} 要么 \(\mathbf{R}_n''\) ”的关系，其中，对于每一个 \(i, \mathbf{R}_i'\) 是这两个关系之一 \(\mathbf{R}_i\) ，``不 \(\mathbf{R}_i\) ''。设 \(S_1, S_2, \ldots, S_p\) 为这些关系。设 \(T_1, T_2, \ldots, T_q\) 是形如“ \(S_{i_{1}}\) 并且 \(S_{i_{2}}\) 并且 \(\ldots\) 并且 \(S_{i_{r}}\) ”的关系，其中 \(i_{1}, i_{2}, \ldots, i_{r}\) 是任意严格递增的指标序列。最后，令 \(T_{0}\) 为关系“ \(R_{1}\) 或（非 \(R_{1}\) ）''，这是C中的一个定理。证明每一个由 \(R_{1}, R_{2}, \ldots, R_{n}\) 逻辑构造的关系在C中等价于且仅等价于以下关系之一 \(T_{0}, T_{1}, \ldots, T_{q}\) （首先证明每个关系 \(R_{i}\) 在C中等价于以下关系之一 \(T_{0}, T_{1}, \ldots, T_{q}\) 。如果R是由 \(R_{1}, R_{2}, \ldots, R_{n}\) 逻辑构造的，请以包含R的基底 \(R_{1}, R_{2}, \ldots, R_{n}\) 为基础，逐步论证一个逻辑构造。最后，利用(d)证明唯一性。）
\item
  设R是C中的一个关系，并令 \(R_{1}, R_{2}, \ldots, R_{n}\) 是其逻辑组成部分。那么，R是C中的一个定理，当且仅当，无论将符号0和1如何分配给 \(R_{1}, R_{2}, \ldots, R_{n}\) ，R 所携带的相应符号为 0。
\end{enumerate}

\begin{enumerate}
\def\labelenumi{\arabic{enumi}.}
\setcounter{enumi}{7}
\tightlist
\item
  设 \(R_{1}, R_{2}, \ldots, R_{n}\) 为理论 G 中逻辑上不可约的关系（习题 6）。为每个关系分配 \(R_{i}\) 0、1、2这些符号之一。对于以基础为基的逻辑构造中的每个关系 \(R_{1}, R_{2}, \ldots, R_{n}\) 我们根据以下符号规则（练习7）分配符号0、1、2中的一个：
\end{enumerate}

\[
\begin{array}{c} \neg 0 = 1, \quad \neg 1 = 0, \quad \neg 2 = 2; \\ \lor 0 0 = \lor 0 1 = \lor 0 2 = \lor 1 0 = \lor 2 0 = \lor 2 2 = 0; \\ \lor 1 1 = 1, \quad \lor 1 2 = \lor 2 1 = 2. \end{array}
\]

\begin{enumerate}
\def\labelenumi{(\alph{enumi})}
\item
  如果 R 是由 \(R_{1}, R_{2}, \ldots, R_{n}\) ，R所携带的符号独立于逻辑构造，以基 \(R_{1}, R_{2}, \ldots, R_{n}\) ，其中R出现。
\item
  假设 \(\mathfrak{G}\) 的公理仅由方案S2、S3、S4提供。设 \(R\) 成為定理 \(\mathfrak{G}\) 。证明无论将符号0、1、2如何分配给 \(R\) 的逻辑分量， \(R\) 所携带的相应符号都是0。另一方面，如果 \(S\) 在逻辑上不可约且符号为2，则 \((S\) 要么 \(S) \Longrightarrow S\) 所携带的符号为2。由此推出， \(\mathfrak{G}\) 不等价于与 \(\mathfrak{G}\) 具有相同符号且其公理由方案S1、S2、S3、S4提供的理论。
\item
  对于仅基于方案S1、S3、S4或方案S1、S2、S4的理论，建立类似的结果。（分别使用以下规则：
\end{enumerate}

\[
\begin{array}{c} \neg 0 = 1, \quad \neg 1 = 0, \quad \neg 2 = 2, \quad \lor 0 0 = \lor 0 1 = \lor 1 0 = \lor 0 2 = \lor 2 0 = 0, \\ \quad \lor 1 1 = 1, \qquad \lor 1 2 = \lor 2 1 = 1, \qquad \lor 2 2 = 1; \\ \quad \neg 0 = 1, \qquad \neg 1 = 2, \qquad \neg 2 = 0, \\ \quad \lor 0 0 = \lor 0 1 = \lor 1 0 = \lor 0 2 = \lor 2 0 = \lor 2 1 = 0, \\ \quad \lor 1 1 = \lor 1 2 = 1, \qquad \lor 2 2 = 2.) \end{array}
\]

\begin{enumerate}
\def\labelenumi{(\alph{enumi})}
\setcounter{enumi}{3}
\tightlist
\item
  对于仅基于方案S1、S2、S3的理论，建立类似的结果。（使用四个符号0、1、2、3及以下规则：
\end{enumerate}

\[
\lnot 0 = 1, \quad \lnot 1 = 0, \quad \lnot 2 = 3, \quad \lnot 3 = 0,
\]

\[
\vee 0 0 = \vee 0 1 = \vee 1 0 = \vee 0 2 = \vee 2 0 = \vee 0 3 = \vee 3 0 = \vee 2 3 = \vee 3 2 = 0,
\]

\[
\vee 1 1 = 1, \quad \vee 1 2 = \vee 2 1 = \vee 2 2 = 2, \quad \vee 1 3 = \vee 3 1 = \vee 3 3 = 3.)
\]
